% =============================================================================
% Project : THE_PHILOSOPHY_OF_VIABILITY
% Book    : Book I — Ontology of Viability
% Book ID : VS.B01
%
% File    : 4_ASSETS/figures/VS.B01.P01.C02.S01.SS04.P001.FG01.invariance_and_transformation.tex
%
% Section / Chapter / Asset:
%   Chapter 02 — Invariance
%   Figure FG01 — Invariance and Transformation
%
% Asset ID:
%   VS.B01.P01.C02.S01.SS04.P001.FG01.invariance_and_transformation
%
% Purpose :
%   To distinguish initial and resulting configurations under
%   transformation and to illustrate the possible coexistence
%   of preservation and alteration with respect to different
%   determinations or relations.
%
% Conceptual Policy:
%   1. The figure does not presuppose that an invariant exists.
%   2. Preservation and alteration are distinct analytical possibilities.
%   3. They may coexist with respect to different specified features.
%   4. Preservation of one feature does not imply that the entire
%      configuration remains unchanged.
%   5. Neither preservation nor alteration establishes numerical
%      identity or non-identity across configurations.
%   6. An invariance claim requires specification of its
%      transformation, domain, respect, and preservation criterion.
%   7. The figure is schematic and does not assert a universal
%      ontological mechanism.
%
% Requires (preamble):
%   \usepackage{tikz}
%   \usetikzlibrary{arrows.meta, calc, positioning, fit}
% =============================================================================

\begin{figure}[htbp]
    \centering

    \begin{tikzpicture}[
        font=\small,
        node distance=0.9cm,
        % --- base box ---
        box/.style={
            draw, rounded corners=3pt, align=center,
            inner sep=5pt, line width=0.7pt
        },
        % --- node roles ---
        config/.style={box, text width=3.2cm, minimum height=1.9cm,
                       fill=blue!8, draw=blue!55!black},
        transformation/.style={box, text width=2.4cm, minimum height=1.5cm,
                       fill=blue!16, draw=blue!55!black},
        outcome/.style={box, text width=5.0cm, minimum height=2.0cm,
                       fill=gray!10, draw=black!60},
        coexist/.style={box, text width=8.0cm, minimum height=1.6cm,
                       fill=green!8, draw=green!40!black},
        chip/.style={box, text width=2.3cm, minimum height=0.95cm,
                     font=\footnotesize, fill=white, draw=black!55},
        % --- frames and labels ---
        vsgroup/.style={draw=black!45, dashed, rounded corners=5pt,
                        inner sep=8pt, line width=0.5pt},
        elabel/.style={font=\footnotesize\itshape, text=black!75},
        condtitle/.style={font=\small\bfseries, align=center},
        % --- arrows ---
        directed/.style={-{Stealth[length=2.5mm]}, line width=0.9pt},
        possible/.style={-{Stealth[length=2.5mm]}, line width=0.8pt, dashed}
    ]

    % ---------------------------------------------------------------
    % Row 1: initial configuration, transformation, resulting configuration
    % ---------------------------------------------------------------
    \node[config] (initial) {
        \textbf{Initial Configuration}\\[3pt]
        Determinations and relations associated with an existent
    };
    \node[transformation, right=of initial] (transform) {
        \textbf{Transformation}\\[3pt]
        Specified change
    };
    \node[config, right=of transform] (result) {
        \textbf{Resulting Configuration}\\[3pt]
        Determinations and relations following transformation
    };

    \draw[directed] (initial) -- (transform);
    \draw[directed] (transform) -- (result);

    \node[vsgroup, fit=(initial)(transform)(result)] (top) {};

    % ---------------------------------------------------------------
    % Row 2: possible preservation / possible alteration
    % ---------------------------------------------------------------
    \node[outcome, anchor=north east] (preservation) at
        ($(top.south)+(-0.4,-1.6)$) {
        \textbf{Possible Preservation}\\[3pt]
        A specified determination or relation remains preserved
        under transformation
    };
    \node[outcome, anchor=north west] (alteration) at
        ($(top.south)+(0.4,-1.6)$) {
        \textbf{Possible Alteration}\\[3pt]
        Another determination or relation undergoes change
        under transformation
    };

    % Shared stem (no arrowhead), then two arrowed branches
    \draw[line width=0.8pt, dashed] (top.south) -- ++(0,-0.8) coordinate (split);
    \draw[possible] (split) -| (preservation.north);
    \draw[possible] (split) -| (alteration.north);

    \node[elabel, anchor=west] at ($(top.south)+(0.15,-0.4)$) {
        Analytical Assessment of the Result
    };

    % ---------------------------------------------------------------
    % Row 3: possible coexistence
    % ---------------------------------------------------------------
    \coordinate (gap) at ($(preservation.south east)!0.5!(alteration.south west)$);

    \node[coexist, below=1.3cm of gap] (coexistence) {
        \textbf{Possible Coexistence}\\[3pt]
        Preservation and alteration may concern different respects
        of the same transformation
    };

    \draw[possible] (preservation.south) -- ++(0,-0.65)
        -| ([xshift=-1.5cm]coexistence.north);
    \draw[possible] (alteration.south) -- ++(0,-0.65)
        -| ([xshift=1.5cm]coexistence.north);

    % ---------------------------------------------------------------
    % Row 4: conditions of an invariance claim (standalone, no arrow)
    % ---------------------------------------------------------------
    \coordinate (condc) at ($(coexistence.south)+(0,-2.1)$);

    \node[chip, anchor=east] (c2) at ($(condc)+(-0.1,0)$) {Domain};
    \node[chip, anchor=west] (c3) at ($(condc)+(0.1,0)$)  {Respect};
    \node[chip, left=0.2cm of c2]  (c1) {Transformation};
    \node[chip, right=0.2cm of c3] (c4) {Preservation Criterion};

    \node[condtitle] (ctitle) at ($(condc)+(0,0.95)$) {
        Conditions of an Invariance Claim
    };
    \node[vsgroup, fit=(ctitle)(c1)(c2)(c3)(c4)] {};

    \end{tikzpicture}

    \caption{Preservation and Alteration as Distinct Possibilities
    Under Transformation}
    \label{fig:invariance-and-transformation}
\end{figure}

% =============================================================================
% END OF FIGURE FG01
% =============================================================================